\section{Forskingsagenda}
\label{sec:agenda-nn}

Teorien genererer sju umiddelbare forskingsprogram, kvart avleidd frå ein spesifikk proposisjon og kvart med falsifiserbare predikasjonar.

\textbf{1. Kryssdomene morforom-replikasjon} (frå proposisjon 1). Teorien predikerer at ein kvar klasse av designa artefaktar vil vise same morforomsstruktur: ikkje-uniform okkupasjon, treregiontpartisjon og kumulativ ekspansjon. Umiddelbare mål inkluderer bygningsfasadar, plakatypografi og programvaregrensesnitt-oppsett.

\textbf{2. Stilkausalitetseksperiment} (frå proposisjon 6). Ingen eksisterande eksperiment isolerer stileksponering som ein uavhengig variabel på formvariasjon. Falsifiseringseksperimentet: gje to grupper av designarar identiske funksjonsbrev og materialbegrensningar men ulik stileksponering. Dersom dei resulterande formene skil seg \textit{berre} langs dimensjonane som svarer til stileksemplara, held teorien. Dersom stileksponering forskyv form langs funksjonelle dimensjonar, krev teorien revisjon.

\textbf{3. Omsorgsdimensjonalitetsmetrikk} (frå proposisjon 5). Teorien definerer omsorg som talet på morforomsaksar agenten aktivt representerer, men gjev ikkje enno ein operasjonell måleprotokoll. Ein tenkjehøgt-protokollanalyse av designarar under designsesjonar kunne gje ein slik metrikk: tel dei distinkte avgrensningsdimensjonane som vert verbaliserte under sesjonen.

\textbf{4. Agentisk designverktøy-prototype} (frå proposisjon 6). Teorien impliserer ein spesifikk arkitektur for maskinassisterte designverktøy: mennesket set dei kognitive lyskjegleparametrane, og maskinen utfører formgrammatikkreglar innanfor den lyskjegla. Ei prototypeimplementering ville teste om kjeglebreidde korrelerer med designkvalitet.

\textbf{5. Ikkje-europeisk uavhengig konvergens} (frå proposisjon 1). Teorien predikerer at landskapsattraktorar er reelle, ikkje kulturelt kontingente. Dersom kinesiske, japanske og vestafrikanske stoltradisjonar konvergerer mot dei same morforomsattraktorane som europeiske tradisjonar, er attraktorstrukturen stadfestar som ein fysisk-ergonomisk realitet.

\textbf{6. Materiell latens-kvantifisering} (frå proposisjon 2). Teorien predikerer eit systematisk etterslep mellom materialtilgang og full formell utnytting. Dette er testbart over fleire materialovergangar: tre til stål (etterslep: omlag 400 år), stål til plast (omlag 50 år), plast til komposittmateriale, kompositt til biomateriale.

\textbf{7. Flerartleg design-formalisme} (frå proposisjon 5). Teoriens agentdefinisjon er substrat-uavhengig, men dei utarbeidde døma avslører ei praktisk utfordring: flerartleg design krev koordinering av agentar der dei kognitive lyskjeglene opererer på ukommensurable tidsskalaer. Ein formell utviding for å romme ukommensurable lyskjegler; med empirisk testing mot faktiske flerartlege designprosjekt som konstruerte våtmarker, permakultur-system eller urban re-vilting-initiativ; ville teste teoriens skalerbarheit utover menneskesentrert design.

Kvart av desse programma er uavhengig publiserbart, empirisk handterleg og direkte avleidd frå teoriens proposisjonar. Saman utgjer dei ein tiårslong forskingsagenda.
